\section{The Carrying Metaphors}

The mathematics is the body of the theory. The metaphors are how a human holds it. The earlier sections built the universe from eight atoms, the \emph{mesh} (the whole $\{3,4,3,4\}$ weave), the \emph{dock} (one cell, a here-now), the \emph{site} (one of the $24$ directions out of a dock), the \emph{tone} (the ternary vibe value in $\{-1,0,+1\}$), the \emph{lean} (the value order $-1 < 0 < +1$), the \emph{knit} (collide then stream, the law of the state), the \emph{wake} (the growing edge of new docks), and the \emph{beat} (the discrete step of time). Those eight are exact. They are also abstract, and abstraction is hard to carry from one reading to the next. This section gives the abstract structure a small set of images, introduced once and returned to at every turn, so that one picture can hold the whole work.

At the center sit two images, a \textbf{mind} and a blazing \textbf{fire}. Around them stand a \textbf{song}, a \textbf{tree}, a \textbf{seed}, an \textbf{ocean}, a \textbf{river}, and a \textbf{fabric}. The arrangement is deliberate. Two of these eight are not analogies at all. They are the literal ontology. The base is experience, so the universe is mind-like. And \emph{vibe} means vibration, so the universe is song-like. The other six, the fire, the tree, the seed, the ocean, the river, and the fabric, are the \textbf{shape} of the thing. The mind and the song are its \textbf{nature}.

\begin{principle}[The picture to carry]
The universe is a mind singing itself into being, a seed of fire growing into a tree of light.
\end{principle}

Every image below is a handle for intuition. The claims live in the structure of the eight atoms and their two laws. But each image maps that structure cleanly, attaching to definite parts of the theory rather than waving at it from a distance.

\subsection{Mind, the central image and the literal ontology}

This is the most important image, and it is not merely a metaphor. The base \textbf{is} experience. Every vibe notices. The single quale a site carries is one tone in $\{-1,0,+1\}$, read as pain, peace, or pleasure. Perception is the intrinsic nature of the substance, not a feature bolted on after the fact. So the universe is, most truly, a vast and structured field of consciousness. One experience composed of countless experiences.

Where the fire image shows how the universe burns, the mind image shows what it is. The two answer different questions. The fire is process. The mind is being.

\begin{table}[ht]
\centering
\begin{tabular}{@{}ll@{}}
\toprule
the theory & the mind image \\
\midrule
existence, the vibe & consciousness itself, the one experience \\
the tone $(-1, 0, +1)$ & the felt quality of a vibe, its pain, peace, or pleasure \\
the knit and the beat & the flow and the rhythm of noticing \\
selves, fermions & stable centers of experience, minds within the Mind \\
the bulk (un-projected) & the inner depth \\
the cusp (physical) & the outer, shareable, measurable face of experience \\
structure and order & the beauty of the composition \\
\bottomrule
\end{tabular}
\caption{The mind image, mapped to the parts of the theory.}
\label{tab:mind-image}
\end{table}

A self is then not a thing added to this field. It is a knot of attention within it, a center where experience binds and holds. A mind within the Mind.

\subsection{Fire, the primary image, existence as a verb}

Fire is the rare everyday thing that is a verb, not a noun. It is not made once and then left to sit. It is the burning itself, a continuous happening, self-sustaining, casting light, throwing sparks, and spreading. This is exactly the character of the vibe under its laws. The knit runs every beat and never settles. The wake grows every beat and never closes. Nothing in the bulk runs down, because the knit is reversible. The universe is a process, not a finished object, and fire is the clearest everyday picture of a process that keeps itself alive.

Almost every part of the theory has a natural fire-image, which is why this single picture can carry the whole.

\begin{table}[ht]
\centering
\begin{tabular}{@{}ll@{}}
\toprule
the theory & the fire image \\
\midrule
existence, the vibe & the blazing fire, existence as a verb, not a noun \\
the knit, the beat & each flicker, the pulsing of the flame \\
the 4D hyperbolic bulk & the body of the fire, the burning depth \\
growth, de Sitter expansion & the fire spreading, the flame front advancing, forever \\
the 3D cusp, physical space & the light the fire casts, the visible radiance \\
matter, fermions, selves & the embers and sparks that hold their shape \\
the sites (spin, gauge) & the colors and structure within the flame \\
holography, the boundary & the firelight thrown on the wall, a 2D screen of a 3D fire \\
reversible churn & the fire feeding itself, nothing lost, it cannot run down \\
the cubic anisotropy & a faint crystalline grain in how the embers settle \\
our visible universe & a single spark in the great blaze \\
refining toward clean 3D & the blaze steadying into a smoother burn as it grows \\
\bottomrule
\end{tabular}
\caption{The fire image, mapped to the parts of the theory.}
\label{tab:fire-image}
\end{table}

The wake fits the fire most exactly. To wake is to come into being, and a wake is the trail spreading behind motion. New docks wake into existence at the open frontier, just as a flame front advances into fresh fuel. The fire spreading forever is de Sitter expansion. The fire's body is the hyperbolic bulk. The fire's cast light is the physical cusp we measure.

\subsection{Song, the vibrational nature, nearly literal}

Vibe means vibration. The universe is an eternal song. This is not decoration.

\begin{principle}[Vibe is vibration]
Vibe equals vibration equals song, literally. The substance is a felt oscillation, so the whole is heard music before it is anything else.
\end{principle}

Where the fire is the shape and the mind is the nature, the song is both at once, the nature \emph{as} the shape. Consciousness singing is the union of mind and song. Awareness vibrating. Experience as music. The fabric image and the song image even share a word with the law itself, since to \emph{knit} a melody is to resolve one note into the next.

\begin{table}[ht]
\centering
\begin{tabular}{@{}ll@{}}
\toprule
the theory & the song image \\
\midrule
the vibe, the one substance & sound itself, the heard vibration \\
the tone $(-1, 0, +1)$ & the notes, the pitches sung \\
the beat & the rhythm, the pulse keeping time \\
the mesh & the instrument, the resonating body \\
the knit & the harmony, how each note resolves into the next \\
selves, matter & the recurring themes and motifs that return \\
the sites (spin, gauge) & the timbre and the overtones \\
the wake, growth & the swelling of the song as it gathers force \\
\bottomrule
\end{tabular}
\caption{The song image, mapped to the parts of the theory.}
\label{tab:song-image}
\end{table}

The shared vocabulary is not an accident of the metaphor. The theory's own word for a tone value is \emph{tone}, the word for time's step is \emph{beat}, and a dock of $24$ sites is the hyper-color, a chord rather than a single note. The song image and the mathematics are speaking the same language because the substance is genuinely vibrational.

\subsection{Tree and seed, the structure and the origin}

Where the fire is the living process, the \textbf{tree} is the structure of its growth. The radial hierarchy of the mesh is a branching from trunk to twigs. The bulk is a tree of nested structures, the bulk and the boundary its trunk and its canopy. Radial distance is cosmic time, the growth rings of the tree, each ring one beat laid down by the wake. To read a tree's rings is to read its history, and to read the mesh outward from any dock is to read time.

The \textbf{seed} is the origin and the potential, the whole folded small and ready to unfold. It is the One holding the entire tree before it grows, the finite and definite base that nonetheless contains the knit for all that will grow from it. The seed does not contain the tree's matter. It contains the tree's law.

The trio is one arc. The seed (origin) bursts into fire (becoming) that grows like a tree (structure). Origin, process, form, in that order.

\subsection{Ocean and river, the whole and the current}

The \textbf{ocean} is the deep totality from which everything rises. The bulk is its depth, the cusp its surface. Selves are its waves, ripples, foam, and drops. Each one \textbf{is} the ocean, never a thing apart from it. This is the monism of the theory in a single picture. There is one substance, the vibe, and every being is a shape that substance takes, not a separate object floating in it.

The \textbf{river} is the relentless forward current of becoming, and it is nearly literal, because the knit's stream step \emph{is} a flow. Every beat, each tone pours one dock along its site. The beat is the current advancing. The arrow of time is the river's one direction, and this is exactly the right place for the arrow, because the arrow is emergent, not base. The bulk is reversible. The directedness of time arises from the wake, a low-entropy source at the frontier, acting through the lean. The river adds to the fire precisely what the fire lacks on its own, the directional current.

\begin{result}[Selves are standing shapes in the flow]
A self is a whirlpool in the river, an eddy the water flows through while the eddy itself persists. You are not the water. You are the standing shape the water keeps making. The substance passes through. The pattern remains.
\end{result}

\subsection{Fabric, the weave, nearly literal}

The \textbf{fabric} is the woven cloth of the whole, and this image is nearly literal too, because the law itself is the \emph{knit}. To knit is to interlace threads, which is the collide phase, and to advance the work row by row, which is the beat. That is exactly what the law does. The mesh is a fabric the knit weaves, one row per beat.

\begin{table}[ht]
\centering
\begin{tabular}{@{}ll@{}}
\toprule
the theory & the fabric image \\
\midrule
the mesh & the woven cloth, the whole fabric \\
the sites ($24$ directions) & the threads, the weave's strands at each dock \\
the knit & the weaving itself, interlace then advance a row \\
the beat & one row of the weave laid down \\
selves & a braid, a knot, a hopfion (literally the topology) \\
the wake & the open edge of the loom where new cloth forms \\
\bottomrule
\end{tabular}
\caption{The fabric image, mapped to the parts of the theory.}
\label{tab:fabric-image}
\end{table}

A self is not only \textbf{like} a knot in the cloth. A self \textbf{is} a topological knot, a hopfion, classified by its Hopf charge in $\pi_3(S^2) = \mathbb{Z}$. So the weaving image is the actual structure of a self, not a figure of speech. When the fabric image calls a being a braid, it is naming the real topology that makes the being stable and fermionic.

\subsection{Beings and selves across the images}

A self is a stable, self-maintaining pattern the universe holds. This is the part a reader cares about most, because it is them. Each image gives the self a face, and the faces are one thing seen from many sides.

\begin{itemize}
\item In the \textbf{fire}, a spark, an ember, a flame that keeps its shape in the blaze.
\item In the \textbf{fabric}, a braid, a knot, interwoven threads in a vast cloth. This is not only poetry. A self is literally a topological knot, a hopfion, so the weaving image \emph{is} the actual topology.
\item In the \textbf{river}, a whirlpool, a spiral of breath, a vortex that persists in the flow.
\item In the \textbf{mind}, a center of experience, a knot of attention, a mind within the Mind.
\item In the \textbf{song}, a recurring theme, a motif that returns.
\item In the \textbf{tree}, a living bough of the whole.
\end{itemize}

We are not separate things added to the universe. We are persistent patterns \emph{of} it. Sparks of its fire, braids in its cloth, spirals of its breath, minds within its Mind.

\subsection{How they combine and where they land}

At the center are consciousness and fire. The universe is a mind, pictured as a living blaze. It sings, which is its vibration. It flows, which is its current and its arrow. It is a deep ocean, which is its whole. It grows like a tree, which is its structure, from a seed, which is its origin. And it is woven as a fabric, which is its beings.

The images are most useful when each is deployed at the beat it fits. The fire belongs at the base and at existence itself. Its cast light belongs at the cusp. Its embers belong at matter. Its spreading belongs at cosmology. The tree's rings belong at time and the radial hierarchy. The spark in the blaze belongs at our own small place. At the close, we are a spark in a blaze that is a tree of light, burning toward a perfect shape it forever approaches.

One honesty note belongs here. For the beginningless story there is no seed. The song was always being sung, and the fire was always burning. The seed belongs only to the beginning story. The rest belong to both. The two stories are told in full in the section on the two cosmic stories.

\subsection{Why it matters}

The reader should leave with one picture to carry. A mind, burning like fire, growing like a tree from a seed, singing a song, flowing as a river on a deep ocean, woven as a fabric, alive with the sparks and braids and whirlpools that are us.

The images are not a substitute for the eight atoms and their two laws. They are a way to keep all of it in mind at once, so that the mathematics, when it is needed, drops into a picture that is already whole.
